\begin{foreignabstract}

This work investigates the application of regression models to predict the performance of a V-AGMD (\textit{Vacuum-Enhanced Air Gap Membrane Distillation}) desalination system, aiming to estimate the permeate flux and the outlet temperatures of the streams from operational process variables. The relevance of the study is associated with the complexity of the heat and mass transfer phenomena involved, which makes it difficult to obtain accurate models using exclusively physical or empirical approaches. Although the permeate flux is adopted as the priority variable in the model selection process, the outlet temperatures are also considered because they provide complementary information about the thermal behavior of the system and are associated with the evaluation of energy indicators, such as the specific energy consumption (SEC) and the gained output ratio (GOR).

The problem is formulated as a multi-output supervised regression task, using experimental data obtained under different operating regimes. Model validation is performed through cross-validation with separation by experimental groups. Different families of models are compared, including regularized linear regressions, tree-based models, and artificial neural networks, with model selection conducted based on the permeate flux prediction error and model simplicity criteria. The results are analyzed in comparison with a simplified physical model, allowing the advantages and limitations of data-driven approaches to be discussed, as well as highlighting the potential of hybrid strategies to improve the generalization capability in process modeling.

\end{foreignabstract}

